\Needspace{6\baselineskip}
\begin{center}\hypertarget{native-250-article178}{}\pdfbookmark[2]{178.}{article-178}178.\end{center}
Functio modo eruta omni quidem rigore errorum probabilitates exprimere
certo non potest: quum enim errores possibiles semper limitibus certis
coërceantur, errorum maiorum probabilitas semper evadere deberet $=0$, dum
formula nostra semper valorem finitum exhibet. Attamen hic defectus, quo omnis
functio analytica natura sua laborare debet, ad omnes usus practicos nullius
momenti est, quum valor functionis nostrae tam rapide decrescat, quamprimum
$h\Delta$ valorem considerabilem acquisivit, ut tuto ipsi $0$ aequivalens censeri
possit. Praeterea ipsos errorum limites absoluto rigore assignare, rei natura
nunquam permittet.

\SourceLeaf{251}Ceterum constans $h$ tamquam mensura praecisionis observationum considerari
poterit. Si enim probabilitas erroris $\Delta$ in aliquo observationum systemate
per $\frac{h}{\BareRoot{\pi}}e^{-hh\Delta\Delta}$, in alio vero systemate
observationum magis minusve exactarum per
$\frac{h'}{\BareRoot{\pi}}e^{-h'h'\Delta\Delta}$ exprimi concipitur,
exspectatio, in observatione aliqua e systemate priori errorem inter limites
$-\delta$ et $+\delta$ contineri, exprimetur per integrale
$\int\frac{h}{\BareRoot{\pi}}e^{-hh\Delta\Delta}\,\mathrm{d}\Delta$,
a $\Delta=-\delta$ usque ad $\Delta=+\delta$ sumtum, et perinde exspectatio,
errorem alicuius observationis e systemate posteriori limites $-\delta'$ et
$+\delta'$ non egredi, exprimetur per integrale
$\int\frac{h'}{\BareRoot{\pi}}e^{-h'h'\Delta\Delta}\,\mathrm{d}\Delta$,
a $\Delta=-\delta'$ usque ad $\Delta=+\delta'$ extensum: ambo autem integralia
manifesto aequalia fiunt, quoties habetur $h\delta=h'\delta'$. Quodsi igitur
e. g. $h'=2h$, aeque facile in systemate priori error duplex committi poterit
ac simplex in posteriori, in quo casu observationibus posterioribus secundum
vulgarem loquendi morem praecisio duplex tribuitur.
